\textbf{Kurzfassung}: Wir alle mussten uns wohl schon einmal mit ihnen
auseinandersetzen, auch wenn die wenigsten von uns gerne darüber
sprechen: Mehlwürmer, Kleidermotten, Ameisen, Mäuse, Silber- und
Papierfischchen, Schaben, Teppichkäfer und ähnliches Getier, das in der
Wohnung umherkreucht und fleucht. Zusammengefasst werden diese Tiere oft
als Schädlinge bezeichnet. Denn auch wenn sie biologisch betrachtet
durchaus wichtige Funktionen wahrnehmen, so richten sie in und an
menschlichen Behausungen doch oft beträchtlichen Schaden an.
Kultureinrichtungen wie Archive, Bibliotheken und Museen sind von dieser
Gefahr nicht ausgenommen und haben aufgrund des in ihnen bewahrten, oft
einzigartigen Kulturguts eine besondere Sorgfaltspflicht auch
hinsichtlich der Schädlingsbekämpfung. Dieser Beitrag soll anhand einer
kleinen Wanderausstellung, die vom Januar bis April 2026 am Deutschen
Museum in München Station machte, einen Überblick über Hintergründe,
Methoden und aktuelle Entwicklungen der sogenannten \enquote{Integrierten
Schädlingsbekämpfung} (oft wird auch im Deutschen der englische
Fachbegriff IPM für \enquote{Integrated Pest Management} verwendet) geben.

\begin{center}\rule{0.5\linewidth}{0.5pt}\end{center}

\noindent\textbf{Abstract}: Most of us have likely had to deal with them at some
point, even if few of us enjoy talking about it: mealworms, clothes
moths, ants, mice, silverfish and paperfish, cockroaches, carpet
beetles, and similar critters crawling and scurrying around the home.
Collectively, these creatures are often referred to as pests. Although
they perform vital biological functions in nature, they often cause
considerable damage within and to human dwellings. Cultural institutions
such as archives, libraries, and museums are not exempt from this threat
and, due to the often unique cultural heritage they preserve, bear a
special responsibility regarding pest control. Based on a small
traveling exhibition that was hosted at the Deutsches Museum in Munich
from January to April 2026, this article provides an overview of the
background, methods, and current developments in the so called
\enquote{Integrierte Schädlingsbekämpfung} (often referred to by the English
term IPM for \enquote{Integrated Pest Management}, even in German).
